% Generated draft: requires mathematical and source review.
\documentclass[11pt]{article}
\usepackage[margin=25mm]{geometry}
\usepackage{fontspec}
\setmainfont{Latin Modern Roman}
\usepackage{amsmath,amssymb,mathtools,unicode-math}
\setmathfont{Latin Modern Math}
\usepackage{enumitem,needspace,xurl,hyperref}
\setlength{\parindent}{0pt}
\setlength{\parskip}{5pt}
\setlength{\emergencystretch}{3em}
\newcommand{\N}{\mathbb{N}}
\newcommand{\Z}{\mathbb{Z}}
\newcommand{\Q}{\mathbb{Q}}
\newcommand{\R}{\mathbb{R}}
\newcommand{\C}{\mathbb{C}}
\newcommand{\F}{\mathbb{F}}
\newcommand{\A}{\mathbb{A}}
\newcommand{\PP}{\mathbb{P}}
\newcommand{\E}{\mathbb{E}}
\newcommand{\eps}{\varepsilon}
\newcommand{\dd}{\,\mathrm{d}}
\begin{document}
\section*{UnsolvedMath ID: 10900016}
\textbf{Source problem code:} catalogue entry
\subsection{Definitions and mathematical statement}
Let $M$ be a $3$-manifold, and let $\phi:\pi_1(M)\to\Z$ be dual to a properly embedded cooriented surface $(\Sigma,\partial\Sigma)\subset(M,\partial M)$. In the customary compact oriented setting this means that the cohomology class of $\phi$ is Poincar\'e--Lefschetz dual to $[\Sigma]\in H_2(M,\partial M;\Z)$. Let $N$ be the manifold obtained by cutting along $\Sigma$, with the two resulting boundary copies $\Sigma_+$ and $\Sigma_-$.

For $\alpha:\pi_1(M)\to\operatorname{SL}_2(\C)$, let $E_\alpha$ be the rank-two complex local system associated with its standard representation, restricted to $N$ and its boundary copies. Say that $N$ is an $\alpha$-twisted homology product if both inclusions induce isomorphisms
\[
 H_i(\Sigma_\pm;E_\alpha)\longrightarrow H_i(N;E_\alpha)
 \quad\text{for every }i,
\]
equivalently $H_i(N,\Sigma_\pm;E_\alpha)=0$ for all $i$ and both signs.

Agol's question, item 3.8, asks to characterize exactly those $\alpha\in\operatorname{Hom}(\pi_1(M),\operatorname{SL}_2(\C))$ for which this holds. The excerpt does not impose connectedness, primitivity or nonvanishing of $\phi$, or norm-minimizing/taut conditions on $\Sigma$; such conditions cannot be silently added.

The source says that the Dunfield--Friedl--Jackson conjecture would imply the property when $M$ is hyperbolic and $\alpha$ is discrete faithful. This is a supplied conditional expectation whose precise hypotheses, especially on the dual surface, are absent here; it is not an unconditional assertion for every surface dual to $\phi$.
\subsection{Short English statement}
For a manifold cut along a surface dual to an integral cohomology class, determine which two-dimensional complex representations make both boundary copies of the surface twisted-homology equivalent to the cut manifold.
\subsection{Sources}
\begin{itemize}
\item Local registry: \url{data/unsolvedmath/problems.json}, permanent UnsolvedMath ID $10900016$, exact source code \texttt{catalogue entry}, attributed to Agol.
\item Delp, Hoffoss and Manning, \emph{Problems In Groups, Geometry, and Three-Manifolds} (2015), subsection 3.8, PDF page 3: \url{https://arxiv.org/pdf/1512.04620#page=3}; abstract record: \url{https://arxiv.org/abs/1512.04620}. Supplied access date: 2026-07-29; source-tex extraction checked against rendered PDF text by the source record.
\item The source invokes the Dunfield--Friedl--Jackson conjecture and the registry mentions twisted-Alexander/homology follow-ups, without supplying exact papers or the full conditional statement.
\end{itemize}
\end{document}
